% ===== PRÉAMBULE CANONIQUE — à coller VERBATIM en tête de CHAQUE .tex =====
\documentclass[11pt,a4paper]{article}
\usepackage[utf8]{inputenc}\usepackage[T1]{fontenc}\usepackage{lmodern}
\usepackage[french]{babel}
\usepackage{amsmath,amssymb,mathtools}\usepackage{icomma}
\usepackage[version=4]{mhchem}          % \ce{...} : équations & formules
\usepackage{siunitx}\sisetup{locale=FR} % \qty{}{}, \si{}, \ang{}
\usepackage{chemfig}                    % \chemfig{...} : structures développées
\usepackage{xcolor}\usepackage{colortbl}
\usepackage{tikz}\usepackage{pgfplots}\pgfplotsset{compat=1.18}
\usepackage[most]{tcolorbox}\usepackage{enumitem}\usepackage{array,booktabs}
\usepackage[a4paper,margin=2.2cm]{geometry}
\usetikzlibrary{arrows.meta,calc,angles,quotes,positioning,patterns,backgrounds,shapes.geometric}
\definecolor{primary}{RGB}{222,49,99}\definecolor{primarydark}{RGB}{150,22,55}
\definecolor{defcol}{RGB}{0,114,178}\definecolor{methcol}{RGB}{52,92,148}
\definecolor{excol}{RGB}{0,135,90}\definecolor{attncol}{RGB}{200,40,55}
\tcbset{boxbase/.style={enhanced,breakable,boxrule=0.9pt,arc=2.5pt,
  left=3mm,right=3mm,top=2.3mm,bottom=2.3mm,fonttitle=\bfseries,coltitle=white}}
% --- boîtes TUTORIEL (noms EXACTS, ne pas renommer) ---
\newtcolorbox{cle}[1][]{boxbase,colback=defcol!6,colframe=defcol,
  title={\textsc{À retenir}\ifx\relax#1\relax\relax\else\textmd{ : \itshape #1}\fi}}
\newtcolorbox{methode}[1][]{boxbase,colback=methcol!7,colframe=methcol,sharp corners,
  title={\textsc{Méthode}\ifx\relax#1\relax\relax\else\textmd{ --- \itshape #1}\fi}}
\newtcolorbox{exemplebox}[1][]{boxbase,colback=excol!5,colframe=excol,
  title={\textsc{Exemple}\ifx\relax#1\relax\relax\else\textmd{ : \itshape #1}\fi}}
\newtcolorbox{piege}[1][]{boxbase,colback=attncol!6,colframe=attncol,
  title={\textsc{Piège}\ifx\relax#1\relax\relax\else\textmd{ : \itshape #1}\fi}}
\newtcolorbox{remarque}[1][]{boxbase,colback=black!4,colframe=black!45,
  title={\textsc{Remarque}\ifx\relax#1\relax\relax\else\textmd{ : \itshape #1}\fi}}
% --- boîtes EXERCICE / CORRIGÉ (noms EXACTS, auto-comptées) ---
\newtcolorbox[auto counter]{exo}[1][]{boxbase,colback=primary!4,colframe=primary,
  title={\textsc{Exercice}~\thetcbcounter\ifx\relax#1\relax\relax\else\textmd{ --- \itshape #1}\fi}}
\newtcolorbox[auto counter]{corrige}[1][]{boxbase,colback=excol!4,colframe=excol,
  title={\textsc{Corrigé}~\thetcbcounter\ifx\relax#1\relax\relax\else\textmd{ --- \itshape #1}\fi}}
\newcommand{\R}{\mathbb{R}}\newcommand{\N}{\mathbb{N}}\newcommand{\Z}{\mathbb{Z}}\newcommand{\ds}{\displaystyle}
% (un \usepackage{fancyhdr}+en-tête est possible ; il est ignoré côté web)
% =========================================================================
\begin{document}
\begin{center}\color{primarydark}\large\bfseries Thème 4 — La formule de Nernst \quad$\bullet$\quad Niveau Moyen\end{center}

\begin{exo}[l'effet de la dilution]
Électrode de cuivre \ce{Cu^{2+} + 2 e^- <=> Cu} ($E_{\mathrm{r}}^{\circ}=+0{,}34$~V, $n=2$). Calcule
son potentiel pour trois concentrations, puis conclus sur l'effet de la dilution.
\begin{enumerate}[label=\alph*)]
  \item $[\ce{Cu^{2+}}]=1$~mol/L
  \item $[\ce{Cu^{2+}}]=0{,}1$~mol/L
  \item $[\ce{Cu^{2+}}]=0{,}01$~mol/L
\end{enumerate}
\end{exo}

\begin{exo}[une demi-pile d'argent]
Électrode d'argent \ce{Ag+ + e^- <=> Ag} ($E_{\mathrm{r}}^{\circ}=+0{,}80$~V, $n=1$) avec
$[\ce{Ag+}]=0{,}10$~mol/L. Calcule son potentiel.
\end{exo}

\begin{exo}[f.é.m.\ hors des conditions standard]
Pile $\ce{Cu | Cu^{2+}(1\,mol/L) || Ag+(0{,}10\,mol/L) | Ag}$ à $25\,^{\circ}$C. Données :
$E_{\mathrm{r}}^{\circ}(\ce{Ag+}/\ce{Ag})=+0{,}80$~V, $E_{\mathrm{r}}^{\circ}(\ce{Cu^{2+}}/\ce{Cu})
=+0{,}34$~V.
\begin{enumerate}[label=\alph*)]
  \item Calcule le potentiel de la cathode (argent).
  \item Calcule le potentiel de l'anode (cuivre).
  \item En déduis la f.é.m.\ $\Delta E$.
\end{enumerate}
\end{exo}

\begin{exo}[le potentiel du permanganate]
Demi-équation \ce{MnO4^- + 8 H+ + 5 e^- <=> Mn^{2+} + 4 H2O} ($E_{\mathrm{r}}^{\circ}=+1{,}51$~V,
$n=5$). On donne $[\ce{MnO4^-}]=[\ce{Mn^{2+}}]=0{,}01$~mol/L et $\mathrm{pH}=0$ (donc $[\ce{H+}]=1$~mol/L).
Le potentiel s'écrit $E=E_{\mathrm{r}}^{\circ}+\dfrac{0{,}06}{5}\log\dfrac{[\ce{MnO4^-}]\,[\ce{H+}]^{8}}
{[\ce{Mn^{2+}}]}$. Calcule $E$.
\end{exo}

\begin{exo}[retrouver un rapport de concentrations]
Pour le couple \ce{Fe^{3+}}/\ce{Fe^{2+}} ($E_{\mathrm{r}}^{\circ}=+0{,}77$~V, $n=1$), on mesure
$E=+0{,}83$~V. Quel est le rapport $\dfrac{[\ce{Fe^{3+}}]}{[\ce{Fe^{2+}}]}$ ?
\end{exo}

\end{document}
